% Part: incompleteness
% Chapter: introduction
% Section: definitions

\documentclass[../../../include/open-logic-section]{subfiles}

\begin{document}

\olfileid{inc}{int}{def}

\olsection{व्याख्या}

गणिताला आकारिक रूप देण्याचा आणि असे रूप सुसंगत व संपूर्ण आहे हे
दाखवण्याचा हिल्बर्टचा कार्यक्रम राबवण्यासाठी प्रथम भाषा, तार्किक चौकट
आणि स्वयंसिद्धक-प्रणाली निवडावी लागेल. आपल्या उद्देशासाठी गणित प्रथम-क्रम
भाषेत आकारिकरीत्या मांडता येते असे समजू. म्हणजे !!{constant}s,
!!{function}s आणि !!{predicate}s यांचा असा काही संच आहे की प्रथम-क्रम
तर्कशास्त्रातील संयोजके व संख्यापक यांच्यासह त्यातून गणितातील विधाने
व्यक्त करता येतात. अशी भाषा अस्तित्वात आहे याबद्दल बहुतेकांचे एकमत
आहे: ती संचसिद्धान्ताची भाषा आहे आणि त्यात $\in$ हे एकमेव अतार्किक
चिन्ह आहे. इतकी साधी भाषा इतकी अभिव्यक्तिक्षम असू शकते हा दावा
पहिल्या नजरेला नक्कीच अविश्वसनीय वाटतो; गणितातील जवळजवळ सर्व काही
या अत्यल्प शब्दसंग्रहात व्यक्त करता येते हे प्रस्थापित करण्यासाठी
मोठे काम करावे लागले. सध्या चर्चा सोपी ठेवण्यासाठी ती केवळ
अंकगणितापुरती, म्हणजे नैसर्गिक संख्या~$\Nat$ यांच्याशी संबंधित
गणिताच्या भागापुरती, मर्यादित करू. अंकगणितातील तथ्ये व्यक्त करण्यासाठी
योग्य भाषा~$\Lang L_A$ आहे. $\Lang L_A$ मध्ये एक द्विस्थानी
!!{predicate}~$<$, एक !!{constant}~$\Obj 0$, एक एकस्थानी
!!{function}~$\prime$, आणि $+$ व $\times$ ही दोन द्विस्थानी
!!{function}s आहेत.

\begin{defn}
!!{sentence}s चा संच~$\Gamma$ हा तार्किक निष्पन्नतेखाली बंद असेल, म्हणजे
$\Gamma = \Setabs{!A}{\Gamma \Entails
!A}$ असेल, तर त्याला \emph{उपपत्ती} म्हणतात.
\end{defn}

उपपत्ती निर्दिष्ट करण्याचे दोन सोपे मार्ग आहेत. एक म्हणजे एखाद्या
!!{structure} मध्ये सत्य असलेल्या !!{sentence}s चा संच म्हणून
तिला निर्दिष्ट करणे. उदाहरणार्थ, $\Lang L_A$ साठी अशी !!{structure}
घ्या की तिचे !!{domain}~$\Nat$ आहे आणि सर्व अतार्किक चिन्हांचा
अपेक्षेप्रमाणे अर्थ लावला आहे.

\begin{defn}\ollabel{def:standard-model}
\emph{अंकगणिताचे प्रमाण प्रतिमान} ही पुढीलप्रमाणे परिभाषित
केलेली !!{structure}~$\Struct
N$ आहे:
\begin{enumerate}
\item $\Domain N = \Nat$
\item $\Assign{\Obj 0}{N} = 0$
\item $\Assign{\Obj \prime}{N}(n) = n + 1$ सर्व $n \in \Nat$ साठी
\item $\Assign{\Obj +}{N}(n, m) = n + m$ सर्व $n, m \in \Nat$ साठी
\item $\Assign{\Obj \times}{N}(n, m) = n\cdot m$ सर्व $n, m \in \Nat$ साठी
\item $\Assign{\Obj <}{N} = \Setabs{\tuple{n, m}}{n \in \Nat, m \in
  \Nat, n < m}$
\end{enumerate}
\end{defn}

`$\times$' आणि `$\cdot$' यांतील फरक लक्षात घ्या: $\times$ हे
अंकगणिताच्या भाषेतील चिन्ह आहे. गुणाकाराची आठवण व्हावी म्हणून ते
निवडले आहे; पण $\times$ ही गुणाकाराची क्रिया नसून द्विस्थानी फलनचिन्ह
आहे (औपचारिकरीत्या, $\Obj f^2_1$). याउलट, $\cdot$ हेच नेहमीचे
गुणाकाराचे फलन \emph{आहे}. $n \cdot m$ सारखे काही दिसल्यास त्याचा
अर्थ $n$ आणि~$m$ या संख्यांचा गुणाकार होतो; $x \times y$ सारखे
काही दिसल्यास आपण अंकगणिताच्या भाषेतील पदाबद्दल बोलत असतो. प्रमाण
प्रतिमानात फलनचिन्ह~$\times$ चा अर्थ नैसर्गिक संख्यांवरील फलन
$\cdot$ असा निश्चित केला जातो. बेरीजेसाठी मात्र अंकगणिताच्या
भाषेतील फलनचिन्ह आणि नैसर्गिक संख्यांवरील बेरीज हे दोन्ही दर्शवण्यासाठी
आपण $+$ वापरतो. त्याचा अर्थ संदर्भावरून ठरवावा लागतो.

\begin{defn}
\emph{सत्य अंकगणिताची} उपपत्ती म्हणजे अंकगणिताच्या प्रमाण प्रतिमानात
पूर्त होणाऱ्या !!{sentence}s चा संच; म्हणजे,
\[
\Th{TA} = \Setabs{!A}{\Sat{N}{!A}}.
\]
\end{defn}

$\Th{TA}$ ही उपपत्ती आहे. कारण $\Th{TA} \Entails !A$ असेल तेव्हा
$\Th{TA}$ ची पूर्ती करणाऱ्या प्रत्येक !!{structure} मध्ये $!A$ ची
पूर्ती होते. $\Sat{N}{\Th{TA}}$ असल्यामुळे $\Sat{N}{!A}$ मिळते;
म्हणून $!A \in \Th{TA}$.

उपपत्ती~$\Gamma$ निर्दिष्ट करण्याचा दुसरा मार्ग म्हणजे काही
वाक्यसंच~$\Gamma_0$ पासून तार्किकरीत्या निष्पन्न होणाऱ्या
!!{sentence}s चा संच म्हणून तिला सांगणे. तेव्हा $\Gamma$ हे
तार्किक निष्पन्नतेखाली $\Gamma_0$ चे ``संवरण'' असते. या मार्गाने उपपत्ती निर्दिष्ट करणे तेव्हाच
उपयुक्त ठरते जेव्हा $\Gamma_0$ स्पष्टपणे निर्दिष्ट केला असेल;
उदाहरणार्थ, $\Gamma_0$ मधील !!{element}s ची यादी दिली असेल.
किमान एवढे तरी हवे की $\Gamma_0$ निर्णेय असावा: एखादे
!!a{sentence} $\Gamma_0$ चे घटक आहे की नाही हे ठरवण्यासाठी
संगणनक्षम चाचणी असावी. $\Gamma_0$ मधील !!{sentence}s ना
$\Gamma$ ची \emph{स्वयंसिद्धके}, आणि $\Gamma$ ला
$\Gamma_0$ या संचातील \emph{स्वयंसिद्धकांद्वारे निरूपित} म्हणतो.

\begin{defn}
उपपत्ती~$\Gamma$ ही $\Gamma_0$ या संचातील
\emph{स्वयंसिद्धकांद्वारे निरूपित} आहे तेव्हाच
\[
\Gamma = \Setabs{!A}{\Gamma_0 \Entails !A}
\]
असते.
\end{defn}

\begin{defn}
पुढील वाक्यांनी स्वयंसिद्धकांद्वारे निरूपित केलेली उपपत्ती
$\Th{Q}$ ``रॉबिन्सनची $\Th{Q}$'' म्हणून ओळखली जाते आणि ती
अंकगणिताची अतिशय साधी उपपत्ती आहे.
\begin{align*}
& \lforall[x][\lforall[y][(\eq[x'][y'] \lif \eq[x][y])]] \tag{$!Q_1$}\\
& \lforall[x][\eq/[\Obj 0][x']] \tag{$!Q_2$}\\
& \lforall[x][(\eq[x][\Obj 0] \lor \lexists[y][\eq[x][y']])] \tag{$!Q_3$}\\
& \lforall[x][\eq[(x + \Obj 0)][x]] \tag{$!Q_4$}\\
& \lforall[x][\lforall[y][\eq[(x + y')][(x + y)']]] \tag{$!Q_5$}\\
& \lforall[x][\eq[(x \times \Obj 0)][\Obj 0]] \tag{$!Q_6$}\\
& \lforall[x][\lforall[y][\eq[(x \times y')][((x \times y) + x)]]] \tag{$!Q_7$}\\
& \lforall[x][\lforall[y][(x < y \liff \lexists[z][\eq[(z' + x)][y]])]] \tag{$!Q_8$}
\end{align*}
$\{!Q_1, \dots, !Q_8\}$ हा !!{sentence}s चा संच $\Th{Q}$ ची
स्वयंसिद्धके आहेत. म्हणून त्यांच्यापासून तार्किकरीत्या निष्पन्न होणारी सर्व
!!{sentence}s $\Th{Q}$ मध्ये आहेत:
\[
\Th{Q} = \Setabs{!A}{\{!Q_1, \dots, !Q_8\} \Entails !A}.
\]
\end{defn}

\begin{defn}
$!A(x)$ हे $\Lang L_A$ मधील असे !!a{formula} आहे असे समजा की
त्यात $x$ आणि $y_1$, \dots, $y_n$ ही मुक्त चरे आहेत. तर पुढील
रूपातील कोणतेही !!{sentence} हे \emph{विगमन-रूपबंधाचे} उदाहरण आहे:
\[
\lforall[y_1][\dots\lforall[y_n][((!A(\Obj 0) \land \lforall[x][(!A(x)
\lif !A(x'))]) \lif \lforall[x][!A(x)])]]
\]

\emph{पेआनो अंकगणित}~$\Th{PA}$ ही $\Th{Q}$ च्या स्वयंसिद्धकांसह
विगमन-रूपबंधाच्या सर्व उदाहरणांनी स्वयंसिद्धकांद्वारे निरूपित
केलेली उपपत्ती आहे.
\end{defn}

\begin{explain}
विगमन-रूपबंधाचे प्रत्येक उदाहरण~$\Struct{N}$ मध्ये सत्य असते.
!!{formula}~$!A$ मध्ये फक्त एकच मुक्त !!{variable}~$x$ असल्यास
हे सर्वात सहज दिसते. तेव्हा $!A(x)$ हे~$\Struct{N}$ मध्ये
$\Nat$ चा उपसंच~$X_{!A}$ परिभाषित करते. $s(x) = n$ असताना
$\Sat{N}{!A(x)}[s]$ होईल अशा सर्व~$n \in \Nat$ चा~$X_{!A}$ हा संच आहे.
विगमन-रूपबंधाचे संबंधित उदाहरण असे आहे:
\[
((!A(\Obj 0) \land \lforall[x][(!A(x) \lif !A(x'))]) \lif 
  \lforall[x][!A(x)]).
\]
याचा पूर्वपक्ष~$\Struct{N}$ मध्ये सत्य असल्यास $0 \in X_{!A}$;
आणि $n \in X_{!A}$ असेल तेव्हा $n+1$ देखील त्यात असते.
$0 \in X_{!A}$ असल्याने $1 \in X_{!A}$ मिळते. पुन्हा
$1 \in X_{!A}$ असल्याने $2 \in X_{!A}$ मिळते, आणि पुढे असेच.
म्हणून प्रत्येक
$n \in \Nat$ साठी $n \in X_{!A}$. याचा अर्थ
$\lforall[x][!A(x)]$ ची~$\Struct{N}$ मध्ये पूर्ती होते.
\end{explain}

$\Th{Q}$ आणि $\Th{PA}$ या दोन्ही स्वयंसिद्धकांद्वारे निरूपित
उपपत्ती आहेत. त्या किती प्रबळ आहेत, हा मुख्य प्रश्न आहे.
उदाहरणार्थ, $\Lang L_A$ मध्ये व्यक्त करता येणारी~$\Nat$ विषयीची
सर्व सत्ये $\Th{PA}$ सिद्ध करू शकते का? नेमके सांगायचे तर,
$\Lang L_A$ मध्ये मांडता येणारे सर्व प्रश्न $\Th{PA}$ ची
स्वयंसिद्धके सोडवतात का?

हा प्रश्न दुसऱ्या रीतीनेही विचारता येईल: $\Th{PA} = \Th{TA}$
आहे का? $\Th{TA}$ अर्थातच~$\Nat$ विषयीची सर्व सत्ये सिद्ध
करते (म्हणजे ती त्यात समाविष्ट असतात), आणि $\Lang L_A$ मधील
सर्व प्रश्न सोडवते. कारण $!A$ हे $\Lang L_A$ मधील !!a{sentence}
असेल तर $\Sat{N}{!A}$ किंवा $\Sat{N}{\lnot !A}$ यांपैकी
एक सत्य असते; म्हणून $\Th{TA}
\Entails !A$ किंवा $\Th{TA} \Entails \lnot !A$. अशा उपपत्तीला
\emph{!!{complete}} म्हणतात.

\begin{defn}
उपपत्ती $\Gamma$ ही \emph{!!{complete}} आहे तेव्हाच तिच्या
भाषेतील प्रत्येक !!{sentence}~$!A$ साठी $\Gamma \Entails !A$
किंवा $\Gamma \Entails \lnot !A$.
\end{defn}

\begin{explain}
संपूर्णता प्रमेयामुळे $\Gamma \Entails !A$ तेव्हाच
$\Gamma \Proves
!A$. म्हणून $\Gamma$ संपूर्ण आहे तेव्हाच तिच्या भाषेतील
प्रत्येक !!{sentence}~$!A$ साठी $\Gamma \Proves !A$ किंवा
$\Gamma \Proves \lnot !A$.
\end{explain}

आणखी एक प्रश्न असा आहे: एखादे !!a{sentence}~$\Th{TA}$ मध्ये,
$\Th{PA}$ मध्ये, किंवा अगदी~$\Th{Q}$ मध्येही आहे की नाही हे
तपासण्यासाठी संगणकीय प्रक्रिया वापरता येईल का? संच (उदाहरणार्थ,
!!{sentence}s चा संच) !!{decidable} कधी असतो हे परिभाषित करून हा
प्रश्न अधिक नेमका करता येईल.

\begin{defn}
संच $X$ हा \emph{!!{decidable}} आहे तेव्हाच अशी संगणकीय प्रक्रिया
असते की आदान~$x$ दिल्यावर $x \in X$ असल्यास ती~$1$ आणि
अन्यथा~$0$ परत करते.
\end{defn}

म्हणून प्रश्न असा होतो: $\Th{TA}$ ($\Th{PA}$, $\Th{Q}$) !!{decidable} आहे का?

या सर्व प्रश्नांचे उत्तर नकारार्थी असेल. यांपैकी कोणतीही उपपत्ती
निर्णेय नाही. मात्र ही घटना केवळ या उपपत्तींनाच लागू नाही. प्रत्यक्षात,
विशिष्ट अटी पूर्ण करणाऱ्या \emph{कोणत्याही} उपपत्तीसाठी हेच निष्कर्ष
लागू होतात. यांपैकी एक अट, जी $\Th{Q}$ आणि $\Th{PA}$ पूर्ण करतात,
अशी आहे की उपपत्ती !!a{decidable} स्वयंसिद्धकसंचाने निरूपित झालेली असावी.

\begin{defn}
एखादी उपपत्ती !!a{decidable} स्वयंसिद्धकसंचाने निरूपित झाली असेल तर ती
\emph{!!{axiomatizable}} आहे.
\end{defn}

\begin{ex}
सांत !!{sentence}s च्या संचाने स्वयंसिद्धकांद्वारे निरूपित
झालेली कोणतीही उपपत्ती !!{axiomatizable} असते, कारण प्रत्येक
सांत संच !!{decidable} असतो. म्हणून, उदाहरणार्थ, $\Th{Q}$ ही
!!{axiomatizable} आहे.

$\Th{PA}$ सारखी रूपबंधांद्वारे स्वयंसिद्धके दिलेली उपपत्तीदेखील
!!{axiomatizable} असते. कारण $!B$ हे $\Th{PA}$ च्या स्वयंसिद्धकांपैकी
आहे का, म्हणजे $!B$ हे $\Th{PA}$ चे स्वयंसिद्धक असल्यास
$\Char{X}(!B) = 1$ आणि
अन्यथा $= 0$ असणारे फलन $\Char{X}$ कसे काढायचे, हे पुढीलप्रमाणे
तपासता येते. प्रथम, $!B$ हे $\Th{Q}$ च्या स्वयंसिद्धकांपैकी आहे
का ते तपासा. असल्यास उत्तर ``होय'' आणि $\Char{X}(!B) = 1$.
नसल्यास ते विगमन-रूपबंधाचे उदाहरण आहे का हे तपासा. हे पद्धतशीरपणे
करता येते. येथे पुढील पद्धत कदाचित सर्वात स्पष्ट आहे: विगमन-रूपबंधाचे
प्रत्येक उदाहरण सुरुवातीला काही वैश्विक संख्यापकांनी, त्यानंतर
सोपाधिक असलेल्या एका उप-!!{formula} ने सुरू होते. त्या सोपाधिकाचा
उत्तरपक्ष $\lforall[x][!A(x, y_1, \dots, y_n)]$ असतो. येथे $x$
आणि $y_1$, \dots, $y_n$ ही~$!A$ ची सर्व मुक्त चरे असतात, आणि
$!B$ च्या सुरुवातीचे संख्यापक $y_1$, \dots,~$y_n$ ही चरे बद्ध
करतात. हे~$!A$ काढून घेतल्यावर त्याची मुक्त चरे सुरुवातीच्या
वैश्विक संख्यापकांनी आणि~$\lforall[x]$ ने बद्ध केलेल्या चरांशी
जुळतात का ते तपासा. मग सोपाधिकाचा पूर्वपक्ष पुढील स्वरूपाशी जुळतो
का ते तपासा:
\[
!A(\Obj 0, y_1, \dots, y_n) \land \lforall[x][(!A(x, y_1, \dots, y_n)
\lif !A(x', y_1, \dots, y_n))]
\]
तो जुळल्यास $!B$ विगमन-रूपबंधाचे उदाहरण आहे; न जुळल्यास $!B$ तसे नाही.
\end{ex}

या प्रश्नाचे---आणि कोणत्या उपपत्ती संपूर्ण किंवा निर्णेय असतात या
अधिक व्यापक प्रश्नाचे---उत्तर शोधताना पुढील व्याख्याही उपयुक्त ठरेल.
संच $X$ हा रिक्त असेल किंवा $f \colon \Nat
\to X$ असे !!a{surjective} फलन असेल तेव्हाच तो !!{enumerable}
असतो, हे आठवा. अशा फलनाला $X$ चे प्रगणन म्हणतात.

\begin{defn}
संच $X$ हा रिक्त असेल किंवा त्याचे संगणनक्षम प्रगणन असेल तेव्हाच
त्याला \emph{!!{computably enumerable}} (संक्षेपात !!{c.e.}) म्हणतात.
\end{defn}

!!{axiomatizability} व्यतिरिक्त, अपूर्णतेची प्रमेये लागू होण्यासाठी
उपपत्तीविषयी आणखी एक अट अशी असेल की ती संगणनक्षम फलने आणि
!!{decidable} संबंध यांच्याबद्दलची मूलभूत तथ्ये सिद्ध करू शकेल इतकी
प्रबळ असावी. ``मूलभूत तथ्ये'' म्हणजे संगणनक्षम फलनांच्या प्रत्येक
आदानासाठी त्यांची मूल्ये काय आहेत हे व्यक्त करणारी !!{sentence}s.
आणि ``पुरेशी प्रबळ'' म्हणजे संबंधित उपपत्ती ही वाक्ये आपली प्रमेये
मानतात. उदाहरण म्हणून बेरीज हे सर्वसाधारण संगणनक्षम फलन घ्या.
$2$ आणि $3$ या आदानांसाठी $+$ चे मूल्य~$5$ आहे, म्हणजे
$2+3 = 5$. $2$ आणि $3$ या आदानांसाठी $+$ चे मूल्य~$5$ आहे
असे सांगणारे अंकगणिताच्या भाषेतील वाक्य आहे:
$(\num{2} + \num{3}) = \num{5}$. उदाहरणार्थ, $\Th{Q}$ हे वाक्य
सिद्ध करते. अधिक सर्वसाधारणपणे, प्रत्येक संगणनक्षम फलन
$f(x_1, x_2)$ साठी~$\Lang{L_A}$ मध्ये असे !!a{formula}~$!A_f(x_1,
x_2, y)$ हवे की $f(n_1, n_2) = m$ असेल तेव्हा $\Th{Q}
\Proves !A_f(\num{n_1}, \num{n_2}, \num{m})$. अशा रीतीने
$\Th{Q}$ हे $n_1$, $n_2$ या आदानांसाठी~$f$ चे मूल्य~$m$
आहे हे सिद्ध करते. प्रत्यक्षात, आपण याहून थोडे अधिक मागतो:
या $n_1$,~$n_2$ या आदानांसाठी~$f$ चे मूल्य दुसरी कोणतीही संख्या नाही हेदेखील
तिने सिद्ध करावे. !!{decidable} संबंधांबाबतही तसेच आहे. पुढील
दोन व्याख्या हे नेमके करतात.

\begin{defn}
!!{formula}~$!A(x_1, \dots, x_k, y)$ हे
$f\colon \Nat^k \to \Nat$ या फलनाचे~$\Gamma$ मध्ये
\emph{!!{represents}} तेव्हाच की $f(n_1,
\dots, n_k) = m$ असेल तेव्हा पुढील दोन्ही अटी पूर्ण होतात:
\begin{enumerate}
\item $\Gamma \Proves !A(\num{n_1}, \dots, \num{n_k}, \num{m})$, आणि
\item $\Gamma \Proves \lforall[y](!A(\num{n_1}, \dots, \num{n_k},
y) \lif y = \num{m})$.
\end{enumerate}
\end{defn}

\begin{defn}
!!{formula}~$!A(x_1, \dots, x_k)$ हे
$R \subseteq \Nat^k$ या संबंधाचे \emph{!!{represents}} तेव्हाच की
\begin{enumerate}
\item $R(n_1, \dots, n_k)$ असेल तेव्हा $\Gamma \Proves !A(\num{n_1},
\dots, \num{n_k})$, आणि
\item $R(n_1, \dots, n_k)$ नसेल तेव्हा $\Gamma \Proves \lnot
!A(\num{n_1}, \dots, \num{n_k})$.
\end{enumerate}
\end{defn}

अपूर्णतेची प्रमेये लागू होण्याइतकी एखादी उपपत्ती ``प्रबळ'' असते, जर
ती सर्व संगणनक्षम फलने आणि सर्व !!{decidable} संबंध !!{represents}
असेल. $\Th{Q}$ आणि तिच्या विस्तारित उपपत्ती ही अट पूर्ण करतात.
पण हे सिद्ध करण्यास आपल्याला थोडा वेळ लागेल: $\Th{Q}$ काय
सिद्ध करू शकते याविषयीचे हे सोपे नसलेले तथ्य आहे, कारण $\Th{Q}$ कडे
थोडीच स्वयंसिद्धके आहेत आणि त्यांच्यापासून ही सर्व तथ्ये सिद्ध
करावी लागतील. तरी $\Th{Q}$ ही अतिशय दुर्बल उपपत्ती आहे.
म्हणून $\Th{Q}$ सर्व संगणनक्षम फलने निरूपित करते हे सिद्ध
करणे कठीण असले, तरी बहुतेक महत्त्वाच्या उपपत्ती~$\Th{Q}$ पेक्षा
प्रबळ असतात, म्हणजे $\Th{Q}$ पेक्षा अधिक गोष्टी सिद्ध करतात.
$\Th{Q}$ ने सिद्ध केलेली गोष्ट कोणतीही अधिक प्रबळ उपपत्तीही
सिद्ध करते; म्हणून $\Th{Q}$ सर्व संगणनक्षम फलने निरूपित करत
असल्यास प्रत्येक अधिक प्रबळ उपपत्तीदेखील तसे करते. त्यामुळे
अपूर्णतेच्या प्रमेयांची ही अट अनेक महत्त्वाच्या उपपत्ती पूर्ण करतात.
आपल्या कठीण परिश्रमांचे फळ म्हणजे ही प्रमेये उपपत्तींच्या विस्तृत
वर्गाला लागू होतात हे दिसेल. नक्कीच, ``सर्व गणित'' आकारिकरीत्या
मांडू पाहणाऱ्या कोणत्याही उपपत्तीने $\Th{Q}$ सिद्ध करते ते
सर्व सिद्ध केले पाहिजे; किमान प्राथमिक संगणनांचे निष्कर्ष तरी
तिला सामावून घेता आले पाहिजेत. म्हणून ``सर्व गणिताची'' उपपत्ती
म्हणवण्यास पात्र असलेल्या कोणत्याही उपपत्तीला अपूर्णतेची प्रमेये
लागू होतील.


\end{document}
